\section{Pruning coherent returns and bicycles}
\label{sec:pruning}

We work with the random universal tables of
Proposition~\ref{prop:sequential-law}.  The coefficient pools have already
been fixed.  The graph in this Section has only old vertices and old edges;
the private completion is made later.  Write $X^{\mathrm{good}}$ for the
graph induced on the vertices surviving the gate and separation deletions
of Proposition~\ref{prop:tower-labels}.  These deletions are defined for
every realization of the tables.  A statement that a path is in
$X^{\mathrm{good}}$ will be used as a condition on the objects being
counted, never as additional conditioning of the table distribution.

All lengths in this Section count edges.  A bicycle is a finite connected
subgraph with first Betti number two and with no vertex of valence one.
After vertices of valence two are suppressed, it is a theta, a
figure-eight, or a barbell.  We need to exclude embedded bicycles; their
edge labels are automatically immersed because the ambient labeling is
folded.

Our goal, stated precisely in Proposition~\ref{prop:pruning}, is to remove
an arbitrarily small fraction of every layer while excluding short cycles
and bicycles. We first count repeated-prefix returns, which yield the
$.93m_i$ girth bound, and then use that bound to count the remaining short
thetas. Removing bicycles below $1.52m_i$ finally ensures that nonempty cyclically
reduced closed paths of length at most $6.04m_i$ are not Heisenberg laws,
the separation input needed for
the voltage choice.

\subsection{Prefix depth, fixed margins, and admissible returns}

We use the aligned depths of \eqref{eq:aligned-depths}, writing
$d_A=\delta_A$ and $d_B=\delta_B$.  Thus, for an A-vertex, agreement to depth
\[
  d_A(j)=(2j+1)s \qquad (j\geq0)
\]
means agreement of both raw streams through length $r+js$.  For a
B-vertex, agreement to depth
\[
  d_B(j)=2js \qquad (j\geq0)
\]
means agreement of both raw streams through length $js$.
Depth is measured in state bits; on the A-side it subtracts the constant
$2r-s$ from the number of specified raw bits.  This is why the common
scale in both cases is $H_i$.

Let $D(v)$ be the largest such depth for which both specified stream
prefixes are active at $v$.  The two possible terminal cutoffs give
\begin{equation}
\label{eq:pruning-available-depth}
 H_i-2s\leq D(v)\leq H_i
 \quad\text{for }v\in\mathscr A_i\sqcup\mathscr B_i.
\end{equation}
For example, if the active cutoff is inside B-block $j$, the B-input is
partial while its evaluation block is full; the largest common full
B-prefix ends at block $j-1$.  Its missing depth is between $s$ and $2s$.
On the other side the input has ended and the next evaluation block is
partial; the missing depth is between zero and $s$.  The case of a final
A-input is the same with the sides reversed.

Two equal state prefixes followed through one common outward label
propagate as
\begin{equation}
\label{eq:pruning-prefix-propagation}
  A_j\longmapsto B_j,\qquad B_j\longmapsto A_{j-1}.
\end{equation}
Here $A_j$ denotes both A-prefixes through $r+js$, and $B_j$ denotes both
B-prefixes through $js$.  Each propagation therefore loses exactly $s$
of aligned depth.  To verify the first assertion, the equal A-evaluations
through $r+js$, after removing their common label head, determine the B-input
through block $j$ by successive inversion.  The B-evaluation through block
$j$ uses only A-input blocks through $j-1$, already known.  The second
assertion follows from the other equation: the B-evaluation through $js$
determines the A-input through block $j-1$, and its A-evaluation uses earlier
B-inputs.  The correction offsets are equal because the label is equal and
their raw gate windows are earlier than both coupled blocks.  Adjacent
cutoffs do not invalidate the stated full prefixes: the cutoff moves by at
most $b<s$, whereas one propagation has removed one whole interleaved
block.  These observations prove~\eqref{eq:pruning-prefix-propagation}
also when one of the original vertices is at its last complete prefix.

We make the margin independent of the particular subpath.  For $k\geq0$
put
\[
 T_k=H_0(11/10)^k,\qquad
 I_k=[5T_k/6,6T_k/5],\qquad M_k=T_k/25.
\]
A vertex belongs to band $k$ when its height scale belongs to $I_k$.
Each height belongs to at most a fixed number of bands, independent of
$H_0$ and $s$, and
\begin{equation}
\label{eq:pruning-margin-range}
 \frac1{30}\leq\frac{M_k}{H_i}\leq\frac{6}{125}<.05
 \quad(H_i\in I_k).
\end{equation}
We shall use the weaker interval $[.03,.05]$ in numerical estimates.
If a finite old-edge subgraph through height $i$ has total length at most
$7m_i$, its scale variation is at most $7bm_i$.  For sufficiently small
$b/s$, choose $k$ with $T_k\leq H_i<T_{k+1}$.  All its vertex scales then
lie strictly inside $I_k$.  Thus every later short-path application fits
one band.  Unlike a disjoint partition into bands, this choice has room
at every endpoint.

\begin{definition}[Coherent return]
\label{def:coherent-return}
In band $k$, a coherent return is a positive-length nonbacktracking path
\[
  v_0,e_1,v_1,\ldots,e_p,v_p
\]
in $X^{\mathrm{good}}$, entirely in that band, whose endpoints are on
the same side and agree through an available aligned depth
\begin{equation}
\label{eq:coherent-threshold}
 h\geq ps/2+M_k,
 \qquad h\leq\min\{D(v_0),D(v_p)\}.
\end{equation}
It is minimal if no proper positive-length contiguous subpath is a
coherent return for the same band and the same $M_k$.
An exact closed nonbacktracking path is included whenever it satisfies
the available-depth condition.
\end{definition}

The endpoint side makes $p$ even.  For a fixed length, side, and band,
we always use the least aligned $h$ satisfying the first inequality
of~\eqref{eq:coherent-threshold}; it differs from its right-hand side by
less than $2s$.  This loses no return and removes $h$ as an extra
enumeration parameter.  Every coherent return contains a minimal one,
by choosing a shortest eligible contiguous subpath.

The threshold is chosen to make the return count smaller than a layer.
At the leading exponential scale, a length-$p$ pattern at initial scale
$H=H_i$ costs $ps$ bits of label choices and $H-h$ free endpoint bits,
while compatibility at the final active blocks imposes $ps/2$ bits of
constraints.  The
inequality $h\geq ps/2+M_k$ therefore leaves an exponent $H-M_k$,
below the layer exponent $H$.  The estimates below retain the label-rate
slack and control the losses from unequal heights, terminal cutoffs,
and the conditioned table law.

\begin{lemma}[The two junction coefficients are different]
\label{lem:coherent-junction}
At the identified endpoints of a minimal coherent return, the two
outward incident coefficient names are different.
\end{lemma}

\begin{proof}
The pools are injective on the full finite collection of family, type,
slot, and permitted extra-value names, rather than only on the slots
at a single vertex.  Hence equality of the two coefficient names gives
the same family and the same slot data.  At an A-endpoint the label head
$h_{\mathrm{lab}}$ is the first $r$ bits of its evaluation.  At a
B-endpoint its suffix is the first $s-r$ evaluation bits and its bucket
is slot data.  Equality of the raw endpoint prefixes therefore gives
the same full outward label.  Extra-dependent coefficients cause no
exception because those extra values are part of their names.

Remove one edge at each end.  By
\eqref{eq:pruning-prefix-propagation}, the new endpoints agree to
depth at least $h-s$, and the new threshold is
$(p-2)s/2+M_k$.  The prefixes are available at the new endpoints by
the same propagation statement.  The remaining subpath is still in
the same band and consists of good vertices.  If $p>2$, this is a
proper coherent return, a contradiction.  If $p=2$, the two edges
entering its middle vertex have the same outward label there, so
foldedness makes the original path an immediate backtrack.  This too
is impossible.
\end{proof}

\subsection{Reconstructing a counted object}

The count separates choices from constraints.  Once the labels and
heights are fixed, two incident evaluations reconstruct the state at
each internal vertex.  A coherent path leaves only raw u-bits beyond
the common endpoint prefix to be chosen; a theta has no free endpoint
bits.  A formal reconstruction need not lie in the prescribed state
spaces.  Its terminal coordinates must satisfy the active cutoffs, and
at a theta's two branch vertices the third incident evaluation must
agree with the first two.  These are the constraints whose probabilities
we estimate.

We record exactly what is fixed in the enumeration.  A path pattern
consists of its length $p$, initial side and height, band, edge labels,
and the choices of next height allowed by those labels.  There are at
most three choices per edge if one does not use the height residues to
determine the choice uniquely.  A theta pattern consists of three
positive branch lengths, the root position in that abstract graph,
the initial side and height, edge labels, and allowed height changes.
All slot names, coefficients, correction extras, and current cutoffs
are fixed by this information.  Patterns with inconsistent labels or
height assignments are rejected before any random table is queried.
We also reject equal coefficient names among the two incident slots
at an internal vertex or the three slots at a branch vertex.
Such a pair specifies the same incident
slot, by the injectivity of the coefficient pools.  At an actual
vertex that slot supports just one edge, so its repetition would
be an immediate backtrack or would identify two supposed branches.
No immersed path or embedded theta is lost.  At the coherent
junction the corresponding exclusion is justified by
Lemma~\ref{lem:coherent-junction}.

There is a number $L_s$ such that the number of choices of an oriented
label together with its allowed next height is at most $L_s$, and,
by increasing $s$ after the finite type data,
\begin{equation}
\label{eq:pruning-label-count}
  L_s\leq2^{1.001s}.
\end{equation}
Indeed the additional factors here are at most a fixed multiple of
the actual alphabet size.  This deliberately loose bound is used
only in counting, not in the later forest norm estimate.

For a coherent path we identify its two endpoints only through depth
$h$.  After that depth, choose all their remaining active \emph{raw}
u-bits freely and fix them as part of the pattern.  Their inactive
u-bits are fixed to zero.  There are no free v-bits.  For a theta there
are no free endpoint bits.

\begin{lemma}[Unique reconstruction and free-coordinate count]
\label{lem:pruning-reconstruction}
Fix a coherent-path pattern with different junction coefficients,
and fix its free endpoint u-bits.  For any realization of the
universal tables there is at most one assignment of vertex states
realizing that pattern as a coherent path.  If its endpoint scales
are $H_-$ and $H_+$, the number of free bits is at most
\begin{equation}
\label{eq:pruning-free-bits}
  (H_-+H_+)/2-h.
\end{equation}
For a theta pattern, there is at most one assignment of states and
there are no free bits.  During either reconstruction all inputs and
keys queried at a source block are determined before that block's
tables are exposed.
\end{lemma}

\begin{proof}
Choose two incident slots at each theta branch vertex and use the two
incident slots at every degree-two vertex.  At the coherent
junction use the two end slots through depth $h$.  At a free end
beyond $h$ use its prescribed u-block and its single incident
evaluation.

First recover every A-block $A_0$.  Its two evaluation heads are
known from the labels.  If their coefficients are $c_1\ne c_2$,
subtraction recovers the raw u-block, and then the raw v-block.
At later blocks the known correction offsets can be moved to the
right-hand sides, so the same two equations have the form
\begin{equation}
\label{eq:pruning-two-evaluations}
 v+c_1u=E_1,\qquad v+c_2u=E_2.
\end{equation}
The difference $c_1-c_2$ is a nonzero field element and is invertible
on the full padded block.  Thus these equations uniquely give
\[
 u=(c_1-c_2)^{-1}(E_1-E_2),\qquad v=E_1-c_1u.
\]
At a free endpoint, $u$ is already prescribed and the one equation
uniquely gives $v$.

The update of B-block $j$ uses the two values of $P_{A,j-1}$ at the
neighboring A-input blocks already recovered.  The subsequent update
of A-block $j$ uses $P_{B,j}$ at the B-input blocks just recovered.
Each key contains raw blocks strictly before its source block;
current correction windows end before both coupled blocks.  All
these data are known before the current table is read.  This proves
the assertion inductively in the individual source-block order
$A_0,B_1,A_1,B_2,A_2,\ldots$.

The induction is also an unambiguous \emph{formal} reconstruction
when a pattern is invalid: solve on full padded blocks, stop after
all blocks that can meet the candidate cutoffs, and then test the
required zeros, labels, and vertex conditions.  An assignment
violating an inactive-coordinate condition need not define a graph
vertex, but it causes no ambiguity in earlier query data.  Every
actual realization is reproduced by the induction, which proves
the claimed upper bounds on realizations.  Ignoring some final
validity conditions will consequently only increase our counts.

It remains to count the free endpoint u-bits.  The active interleaving
implies
\[
 r-s\leq \ell_A(i)-\ell_B(i)\leq r,
 \qquad H_i=\ell_A(i)+\ell_B(i)+s-r.
\]
Consequently
\[
 \ell_A(i)\leq H_i/2+(2r-s)/2,
 \qquad \ell_B(i)\leq H_i/2.
\]
At aligned depth $h$ an A-endpoint has specified
$(h+2r-s)/2$ u-bits, and a B-endpoint has specified $h/2$ u-bits.
Adding the two residual lengths on their common side gives
\eqref{eq:pruning-free-bits}.  A theta has two chosen evaluations
at every vertex and hence has no freely prescribed coordinate.
\end{proof}

\subsection{The probability estimate at one source block}

Let $\mathcal F_{<t}$ contain all universal tables at source blocks
strictly preceding $t$, in every family, and all fixed pattern data.
The last assertion of Lemma~\ref{lem:pruning-reconstruction} says
that keys, input arguments, gates, coefficient choices, and the
widths of current tests are $\mathcal F_{<t}$-measurable.  In
particular, no v-coordinate solved from a current evaluation is
used in a current raw-u key.

By Proposition~\ref{prop:sequential-law}, an event $B_t$ involving
$k_t$ whole current permutation tables satisfies
\begin{equation}
\label{eq:pruning-current-domination}
 \mathbb P(B_t\mid\mathcal F_{<t})
 \leq \exp(\epsilon_s k_t)
       \mathbb P_{\mathrm{prod}}(B_t\mid\mathcal F_{<t}),
 \qquad \epsilon_s\longrightarrow0.
\end{equation}
The right-hand law is the product of the indicated current uniform
permutation laws, with earlier tables held fixed.  It is legitimate
to apply this bound to a union of output prescriptions, hence to a
linear compatibility condition.  Untested entries of a table in
$B_t$ are integrated out.  The estimate is not an assertion about
the last unknown entry of a permutation after its other entries
have been prescribed.

\begin{lemma}[Predictable accumulated rank]
\label{lem:pruning-rank-process}
Suppose that, at each of finitely many source blocks, a test is
selected from $\mathcal F_{<t}$, and its product-law probability is
at most $2^{-R_t}$.  The nonnegative integers $R_t$ and $k_t$ are
$\mathcal F_{<t}$-measurable.  If an event $A$ implies that all tests
pass, that $\sum_tR_t\geq R$, and that $\sum_tk_t\leq K$, then
\[
 \mathbb P(A)\leq \exp(\epsilon_sK)2^{-R}.
\]
The event $A$ may additionally require later minimality, missed-gate
counts, and full consistency of the reconstructed vertex states.
\end{lemma}

\begin{proof}
After a failed test give the process value zero.  Otherwise multiply
its value at step $t$ by $2^{R_t}\exp(-\epsilon_s k_t)$ when that
test passes.  Conditional expectation does not increase, by
\eqref{eq:pruning-current-domination}.  The initial value is one,
so the final nonnegative variable has expectation at most one.
On $A$ its value is at least $2^R\exp(-\epsilon_sK)$, proving the
claim.  Additional final conditions are used only in this last
lower bound, not in any conditional probability.
\end{proof}

\subsection{Terminal dimensions and equality of keys}

The two-evaluation reconstruction solves for a full padded block, but
only its active coordinates may be nonzero.  At a terminal block this
restriction imposes a number of compatibility conditions determined by
the excess of evaluation bits over input bits.  We first count these
conditions, then determine when different vertices test distinct tables.

For a vertex put
\[
 \lambda_A(i)=r+\ell_B(i)-\ell_A(i),\qquad
 \lambda_B(i)=s-r+\ell_A(i)-\ell_B(i).
\]
These are its surplus evaluation lengths $|v|-|u|$.  They lie in
$[0,s]$, and
\begin{equation}
\label{eq:pruning-surplus-pair}
 \lambda_A(i)+\lambda_B(i)=s,
 \qquad
 |\lambda_A(i)+\lambda_B(j)-s|\leq b
 \quad(|i-j|\leq1).
\end{equation}
The last inequality follows because a height step changes just one
input length by $b$.

\begin{lemma}[Terminal compatibility]
\label{lem:pruning-terminal-codimension}
If both incident common cutoffs equal the receiving vertex's own
cutoff, the terminal constraints from two independent incident
evaluations have product-law probability exactly $2^{-\lambda(v)}$.
If each incident common cutoff is at most $b$ bits below that
receiver cutoff, one can test an event on those two current tables
with probability at most
\[
 2^{-\max\{0,\lambda(v)-4b\}}.
\]
The keys must be distinct between the tested edge occurrences;
the input arguments must be known before the current tables are
exposed, and the required terminal gates must be open.
Only the actual uniform active prefix of a partial permutation is
used in this statement.
\end{lemma}

\begin{proof}
Apart from the initial short block, which is never terminal, there
are two terminal forms.  If the final active input has $t$ bits in
its $s$-bit field, then its evaluation block has all $s$ bits.
The map in~\eqref{eq:pruning-two-evaluations} sends an $s$-bit v
and a $t$-bit u-subspace injectively into two $s$-bit evaluations.
Its image has dimension $s+t$, so the compatibility codimension is
$s-t$.  The nonzero coefficient difference proves injectivity on
the full field and therefore on the active subspace.

On the opposite side, the input has ended and the last evaluation
has $t$ bits.  Its two incoming values are independent uniform
$t$-bit values from the two active-prefix permutations.  Since
$u=0$, they must agree after their known offsets are removed.
This has codimension $t$.  These two surpluses add to $s$.
The boundary cases $t=0,s$ simply move the one nonzero surplus to
the other side.  This proves the equal-cutoff statement.

For adjacent incidences the correction extras are fixed label
data, and the lagged gates are determined by earlier raw blocks.
For incidence $a$ write its known correction as
$o_a=C_a^v+K_aC_a^u$ and move the whole vector $o_a$ to the other
side of the equation.  Its support may fill the output field;
translation by it costs no random dimension.  On a common active
random prefix, with coordinate projection $\pi$ and width $w'$, the
map
\[
 (u,v)\longmapsto
 \bigl(\pi(v+K_1u-o_1),\pi(v+K_2u-o_2)\bigr)
\]
has affine image of dimension at most $w'+t$ in $2w'$ bits.
Thus its codimension is at least $\max\{0,w'-t\}$.  This uses no
invertibility of a truncated multiplier.  For the terminal form
with no input, replace $t$ in this dimension count by zero.

Here are the active widths, including the boundary between blocks.
Write $t\in\{b,2b,\ldots,s\}$ for the width of the last active
block in the interleaved raw input stream at the receiver's height;
each incident common cutoff is the receiver cutoff or is $b$ bits
earlier.  If $t>b$, both evaluations of the
partial-input side occur in the preceding source block and have
the full $s$ random bits.  Their codimension is at least $s-t$.
On the evaluation-only side the common random prefix has width
at least $t-b$, giving that many compatibility bits.  If $t=b$,
the cutoff can cross a block boundary.  On the partial-input side
the preceding source block still has full random width: it is
either ordinary or the full active prefix of a marked block.
It therefore still contributes $s-b$ compatibility bits.  On the
opposite side a lower incidence can already be in its deterministic
post-mark block.  Its surplus is only $b$, so we use the certain
event and charge zero, as permitted by
$\max\{0,b-4b\}=0$.  The same reasoning applies with A and B
interchanged; $t=s$ includes the boundary represented by $t=0$
on the other side.  All retained positive-width outputs are
actual independent uniform current-table prefixes.  These bounds
are stronger than the stated conservative loss of $4b$, and prove
the probability assertion.
\end{proof}

\begin{lemma}[A repeated terminal key identifies receiver prefixes]
\label{lem:pruning-key-prefix}
Suppose two distinct tested terminal edge occurrences at the same
source block use the same table key.  Their receiving vertices
have the same side and agree in raw u and v through available
aligned depth at least
\[
  H_*-8s,
\]
where $H_*$ is the minimum scale among the vertices of the counted
path or theta.  At any nonterminal update, equality of keys
likewise identifies both receiver prefixes through the preceding
complete receiver block.
\end{lemma}

\begin{proof}
At a B-block $j$ update the queried table is $P_{A,j-1}$.  Equality
of its key identifies the family, full label, A raw inputs through
block $j-2$, and B raw inputs through block $j-1$.  In particular
the B u-prefixes through $j-1$ agree.  To compute their B
v-prefixes through $j-1$, the equations use A inputs through
$j-2$, the same earlier keys, and the same coefficient and label
offsets.  Correction gates use still earlier raw pairs.  Thus the
evaluation prefixes agree, and $v=Q-Ku$ gives agreement of the v
prefixes.  At an A-block $j$ update the queried table $P_{B,j}$
identifies A raw inputs through $j-1$ and B raw inputs through
$j-1$.  The A evaluation through $j-1$ uses only this B data and
earlier keys, giving the identical conclusion.

More explicitly, if the current table has source index $k$, its
preceding complete receiver prefix has aligned depth exactly $ks$:
for $P_{A,j-1}$ this is $2(j-1)s$, and for $P_{B,j}$ it is
$(2j-1)s$.  Write $\kappa$ for the final common source index of
an incident edge and $t\in\{b,2b,\ldots,s\}$ for its active
length.  Its common scale is $H_c=\kappa s+t$.  The two
terminal forms test source index $\kappa-1$ or $\kappa$,
respectively.  Thus $ks\geq H_c-2s$.  The common scale is at
least the receiving vertex scale minus $b$, and hence
$ks\geq H_*-2s-b$.  If two adjacent incidences straddle a
block boundary, restricting both to their earlier common
receiver prefix loses at most one more aligned step, of length
$2s$.  This is still greater than $H_*-8s$ for $b<s$.
Every prefix just used ends before the first partial receiving
block and is therefore active at both receivers.
Repeated keys at the same receiving vertex cannot occur among
its two chosen slots: their full label names are different.
\end{proof}

\subsection{A deterministic end guard for minimal returns}

Terminal compatibility probabilities multiply when the tested tables
are distinct.  For a minimal coherent return, a repeated key would
identify a shorter coherent subpath unless both receivers lie near
opposite ends.  We omit terminal tests in fixed end intervals to remove
this exception without conditioning the law on minimality.

Fix a coherent pattern of length $p$, initial scale $H$, and margin
$M=M_k$.  Let $H_*$ be the minimum of its vertex scales.  The
height assignment is fixed and
\begin{equation}
\label{eq:pruning-scale-window}
 H-bp\leq H_*\leq H_v\leq H+bp
 \quad\text{at every candidate vertex}.
\end{equation}
Set
\[
 E_p=100(s+bp),\qquad
 J_p=\left\lceil4E_p/s\right\rceil+4.
\]
The end guard consists of all receiver positions at distance at
most $J_p$ from either end of the abstract path.  It depends only
on the fixed pattern data, not on the tables.

\begin{lemma}[Key separation outside the end guard]
\label{lem:pruning-end-guard}
In every minimal coherent return, all terminal keys tested outside
the end guard are distinct at each source block.
\end{lemma}

\begin{proof}
By the available endpoint depth,
\[
  p\leq2(H_*-M+E_p)/s.
\]
Suppose two keys at terminal receiver positions $a<b$ coincide.
Lemma~\ref{lem:pruning-key-prefix} gives agreement through depth
at least $H_*-E_p$.  The receivers have the same side.  Their
contiguous path segment is nonbacktracking and lies in the same
band.  If
\[
 (b-a)s/2+M\leq H_*-E_p,
\]
it is a proper coherent return, contrary to minimality, unless it
is the entire original path.  The latter case already has both
positions in the guards.  Otherwise
\[
 b-a>2(H_*-M-E_p)/s\geq p-4E_p/s.
\]
It follows that $a<4E_p/s$ and $p-b<4E_p/s$.  Both positions
again lie in the guards.  This proves the claim using only contiguous
subpaths of the original return; the prefix identification does not
close a new path.
\end{proof}

In a counting experiment we reject a pattern at an update if its
selected terminal keys outside the guard coincide, or if a
required terminal gate is closed.  These are predictable
rejections because keys and gates are known in
$\mathcal F_{<t}$.  Lemma~\ref{lem:pruning-end-guard} says that
no actual minimal return in $X^{\mathrm{good}}$ is lost by this
rule.  We do not condition on the return being minimal.

Choose one terminal compatibility test at each unguarded internal
vertex with positive surplus.  At a given update these tests
use distinct tables and their product-law probabilities
multiply.  Pair successive internal vertices along the path in
\eqref{eq:pruning-surplus-pair}.  Omitting the two endpoints and
the guards loses at most $(2J_p+4)s$ surplus bits; the adjacent
cutoff allowance loses at most $4bp$.  Therefore the predictable
terminal tests accumulated by an actual minimal return have
total rank at least
\begin{equation}
\label{eq:pruning-path-rank}
 R_{\mathrm{ret}}
 \geq ps/2-10^4(s+bp).
\end{equation}
The large absolute constant merely dominates the displayed guard,
pairing, and $4b$ losses.  If the right-hand side is negative the
assertion follows from nonnegativity of rank.  At most $2p$ whole
tables are charged.  Untested current outputs still participate
in later formal reconstruction, but are integrated out in the
current estimate.

\begin{lemma}[Expected coherent-return count]
\label{lem:pruning-return-count}
For sufficiently large $s/b$ and then $H_0/s$, the expected number
of minimal coherent returns based on either side at height $i$,
counted over all bands containing them, is at most
\[
 2^{.975H_i}.
\]
This counts oriented rooted paths, including repetitions of
vertices and edges.
\end{lemma}

\begin{proof}
Put $H=H_i$.  The first endpoint has scale $H$, so
\eqref{eq:coherent-threshold} gives $p\leq2(H-M)/s\leq2H/s$.
There are a fixed number of eligible bands and at most $L_s^p$
label/height patterns of a given length.  By
Lemma~\ref{lem:pruning-reconstruction}, their free endpoint
coordinates have at most
\[
 2^{H-h+bp}
\]
choices.  Each pattern and coordinate choice has at most one
actual realization.  Apply Lemma~\ref{lem:pruning-rank-process}
using~\eqref{eq:pruning-path-rank}; the event being bounded is
that this realization is a minimal return in the good graph.
Its probability is at most
\[
 \exp(2\epsilon_s p)
 2^{-ps/2+10^4(s+bp)}.
\]
Thus its contribution, summed over label patterns and free bits,
is at most
\[
 2^{H+.001sp-M+10^5(s+bp)}.
\]
Here we have used $h\geq ps/2+M$, and the larger error constant
also absorbs the exponential inflation; for large $s$,
$\epsilon_s\leq1$, and $b\geq1$.

Since $M\geq.03H$ and $sp\leq2H$, the exponent is at most
\[
 .972H+10^5(s+2bH/s).
\]
Choose, for example, $s/b\geq10^{12}$ and $H_0/s\geq10^{12}$,
in addition to the earlier requirements.  The error here is
less than $.001H$.  The finite sum over lengths and bands is
bounded by a fixed polynomial in $H$, which is absorbed by
another $.002H$ on increasing $H_0$ if necessary.  This proves
the stated bound.  These numerical choices are only convenient
sufficient choices; no degree grows with tower height.
\end{proof}

Delete every vertex on every minimal coherent return in every
band, and denote the resulting induced graph by $X^{(1)}$.

\begin{lemma}[The first deletion gives base girth]
\label{lem:pruning-base-girth}
No nonempty cyclically reduced closed path based at height $i$ in
$X^{(1)}$ has length at most $.93m_i$.
\end{lemma}

\begin{proof}
Such a path lies in one of the interior bands chosen above.  Its
endpoints are the same vertex, and its available prefix depth is
at least $H_i-2s$.  Its required depth is at most
\[
 \left(\frac{.93}{2(.50005)}+.05\right)H_i
 <.98H_i.
\]
The least aligned depth above this number adds less than $2s$.
For sufficiently large $H_0/s$ it is still available.  Hence this
path is a coherent return for that band's margin and contains
a minimal return.  At least one of its vertices, in fact all
vertices of that minimal return, was deleted.  This is a
contradiction.
\end{proof}

\subsection{Theta constraints}

Fix a bicycle of length $e<1.52m_i$ with a root at height $i$ in
$X^{(1)}$, and put $H=H_i$.  Along it
$|H_v-H|\leq be$.  A figure-eight or barbell contains two
edge-disjoint cycles.  By Lemma~\ref{lem:pruning-base-girth},
their total length is greater than
\[
 1.86(1-1.52a)m_i>1.52m_i,
 \qquad a=b/\log_2q,
\]
for our choice of $a$.  Thus only a theta is possible.  Its
shortest branch has length at most $e/3<.51m_i$, and any two
vertices are joined by a path of length at most $e/2<.76m_i$:
two edge-disjoint connecting paths have combined length at
most $e$.  Every theta vertex lies on a cycle of length at
most $e$, so we may also restrict the count to
$e>.9m_i$.

For a theta, the terminal tests will impose approximately $se/2$
compatibility bits.  Against approximately $se$ label bits, this alone
does not put the count below the layer size $2^H$ throughout the required
range of $e$.  The additional saving
comes from the third evaluation at each of the two branch vertices.
We estimate those relations first and then arrange their tests at
different source updates from the terminal tests.

\begin{lemma}[Terminal keys in a surviving theta]
\label{lem:pruning-theta-keys}
Let a theta in $X^{(1)}$ have $e<1.52m_i$ edges and a root at
height $i$, and set $H=H_i$. At every source-block update, the terminal
keys used at its distinct vertices are distinct. If its two branch
vertices have the same side, any key shared by their branch
triples at a depth beyond $.65H$ is impossible, apart from a
fixed initial/final block allowance that is omitted below.
\end{lemma}

\begin{proof}
A repeated terminal key gives receiver-prefix agreement through
$H-be-8s$ by Lemma~\ref{lem:pruning-key-prefix}.  Join the
receivers by a shortest theta path.  It has length below $.76m_i$
and lies in an interior band containing the theta.  Its coherent
threshold, including its margin, is at most
\[
 \left(\frac{.76}{2(.50005)}+.05\right)H
 <.81H.
\]
This is smaller than the available agreement depth for sufficiently
small $b/s$ and large $H_0/s$.  The path is a coherent return,
contradicting its survival in $X^{(1)}$.

For a shared key at the branch vertices, the last sentence of
Lemma~\ref{lem:pruning-key-prefix} identifies their preceding
complete receiver prefixes.  A source update at aligned depth
at least $.65H$ identifies them to depth at least $.65H-4s$.
Their shortest branch has length below $.51m_i$, whose coherent
threshold is less than
\[
 \left(\frac{.51}{2(.50005)}+.05\right)H<.56H.
\]
Again the available margin is much larger than the block and
height errors.  This gives the same contradiction.  Equality
of unordered triples in particular would give a shared key.
\end{proof}

At a full ordinary receiving block, the three evaluations at a branch vertex
satisfy
\begin{equation}
\label{eq:pruning-triple-relation}
 (c_2-c_3)E_1+(c_3-c_1)E_2+(c_1-c_2)E_3=0.
\end{equation}
The three coefficients are nonzero.  Corrections have been
incorporated into the known offsets of the $E_j$.

\begin{lemma}[Ranks at the two branch vertices]
\label{lem:pruning-triple-ranks}
One branch-vertex relation in an ordinary full block with three open
gates has product-law probability $2^{-s}$.  If two same-side
branch vertices have different unordered key triples, their two
relations have joint probability $2^{-2s}$.  If their triples
are the same, it is always safe to charge just one relation.
\end{lemma}

\begin{proof}
The three slots at one branch vertex have different labels and hence
different keys.  At their already known input arguments, the
three fresh table values are independent uniform field
elements.  Equation~\eqref{eq:pruning-triple-relation} can be
solved for any one of them and has probability $2^{-s}$.

Two distinct three-element key sets each contain a key absent
from the other set.  Condition on all the \emph{other whole
tables}, leaving these two exclusive tables unexposed.  Each
exclusive table is used at one known argument in its branch-vertex
equation.  Its value is uniform and has a nonzero coefficient;
it does not enter the other branch equation.  The two equations
therefore have probability $2^{-2s}$.  If other formal recovery
queries use those same tables, their results have not been
prescribed and are integrated out.  With identical triples
retain only one of the two necessary equations, giving the
last assertion.
\end{proof}

We now separate branch tests from terminal tests explicitly.
All terminal updates of the fixed theta occur in the source-index
interval
\[
 \left[\left\lfloor\frac{H-be}{s}\right\rfloor-4,
       \left\lceil\frac{H+be}{s}\right\rceil+2\right].
\]
Indeed its common incidence scales lie between $H-be-b$ and
$H+be$, and the calculation in
Lemma~\ref{lem:pruning-key-prefix} places every terminal test
at $\kappa-1$ or $\kappa$ when $H_c=\kappa s+t$.
The displayed interval contains at most $9+2be/s$ integers.
Enlarge this interval by ten
indices at either end and omit it from branch charges.  Also
omit the first ten source indices and every marked split block
or deterministic post-mark identity block at a branch-vertex slot.
These omissions are fixed by the pattern and the universal
mark locations.  In particular, a current terminal test and a
current branch test never use the same source-block update.

The mark spacing has been chosen so that, over the finite union
of the six branch-vertex slot sequences and separately for both source
parities, the deterministic omitted blocks cost less than
$.001H+O(s+be)$ bits.  The gate deletion is likewise normalized
separately on A-source and B-source blocks of each slot.  A
vertex is good only when the proportion of closed ordinary
gates in each such sequence is below $.005$.  Thus the union
of bad gates at the three slots of each of the two branch vertices
costs at most $.015H+O(s+be)$ bits.  This denominator convention
is needed: a bound normalized on all interleaved indices
without separating parity would permit twice this loss.

Before depth $.65H$, charge only the first chosen branch vertex when
the two branch vertices have opposite sides.  When they have the same
side and their triples coincide, again charge one; when their
triples differ we may charge either one or both.  After that
depth charge both branch vertices when their gates are open, rejecting
any forbidden shared key predictably.  All selection decisions
are made from $\mathcal F_{<t}$.  Each branch vertex has
$H/(2s)+O(1+be/s)$ full updates before the terminal band.  Charging
one before $.65H$ and two afterward gives at least
\[
  (.65/2+.35)H-O(s+be)=.675H-O(s+be)
\]
bits.  Lemmas~\ref{lem:pruning-theta-keys} and
\ref{lem:pruning-triple-ranks}, followed by the two omitted-block
budgets, consequently give
\begin{equation}
\label{eq:pruning-branch-budget}
 R_{\mathrm{branch}}
 \geq .675H-.015H-.001H-10^4(s+be)>.64H.
\end{equation}
The last inequality follows from the same sufficient choices
$s/b\geq10^{12}$ and $H_0/s\geq10^{12}$, enlarged if earlier
finite constants require it.

For terminal tests, choose two slots at every theta vertex.
Lemma~\ref{lem:pruning-theta-keys} gives distinct keys, and the
last-three-update gate requirement makes the necessary gates
open.  Pair adjacent internal vertices on each of its three
branches in~\eqref{eq:pruning-surplus-pair}; there are at most
six unpaired or branch vertices.  The theta has $e-1$ vertices.
Lemma~\ref{lem:pruning-terminal-codimension} therefore gives
\begin{equation}
\label{eq:pruning-theta-terminal}
 R_{\mathrm{term}}
 \geq se/2-100(s+be).
\end{equation}
Every terminal update was omitted from the branch charges.  The two
rank contributions therefore compose chronologically under
Lemma~\ref{lem:pruning-rank-process}.  Exclusivity among branch triples
was used only to compute the branch rank at a single update; it does
not assert independence from terminal tests.

\begin{lemma}[Expected theta count]
\label{lem:pruning-theta-count}
For sufficiently large allowed parameters, the expected number
of rooted embedded bicycles in $X^{(1)}$ of length below
$1.52m_i$, based on either side at height $i$, is at most
\[
 2^{.89H_i}.
\]
\end{lemma}

\begin{proof}
Only theta patterns need be considered.  For a fixed total
length $e$, choose its three ordered positive branch lengths,
the position and orientation of its root, and its bipartite
side.  The number of these choices is at most $100(e+1)^4$.
For example there are fewer than $(e+1)^2$ length triples, at
most $2e$ oriented root positions, two side choices, and a
fixed finite number of choices of the first branch.  Give
its edges orientations for enumeration and choose at most
$L_s^e$ label/height patterns.  Assigning heights along a
spanning tree determines the remaining heights; any
inconsistency on a remaining edge causes rejection.

Lemma~\ref{lem:pruning-reconstruction} has no free coordinates
for a theta.  For each fixed pattern use the predictable
terminal and branch tests above.  Every actual theta in
$X^{(1)}$ passes their key and gate rejections and has the
stated accumulated ranks.  The current-table count is at
most $2(e-1)$ for the terminal tests plus six per branch
update.  Since $e>.9m_i$, the latter count is $O(H_i/s)=O(e)$;
a bound of $20e$ suffices after omitting the fixed initial
updates.  Lemma~\ref{lem:pruning-rank-process} bounds its
probability by
\[
 \exp(20\epsilon_s e)
 2^{-se/2-.64H_i+100(s+be)}.
\]
After summing patterns, the exponent is at most
\[
  (1.001-.5)se-.64H_i+10^5(s+be).
\]
For $e<1.52H_i/(.50005s)$ its main coefficient of $H_i$ is
\[
 \frac{.501(1.52)}{.50005}-.64
 =.88288771122\ldots<.883.
\]
The error and the polynomial sum over $e$ fit below $.007H_i$
by the allowed parameter choices.  This proves the bound.
Embeddedness and survival in $X^{(1)}$ are final predicates on
the reconstructed assignment; they were not used to
condition any current table law.
\end{proof}

Delete every vertex on each bicycle with length less than
$1.52m_i$ for at least one of its vertices of height $i$.
Let $X_{\mathrm{old}}$ be the resulting induced graph.

\subsection{From bicycle pruning to Heisenberg nonlaws}

The counting estimates for both deletions are now established.  Before
choosing a realization with small losses in every layer, we derive the
deterministic consequence of bicycle exclusion that the voltage choice
requires.

We use the convention $[a,b]=aba^{-1}b^{-1}$ and write
\[
  \HH
  =\langle x,y,z\mid [x,y]=z,\ [x,z]=[y,z]=1\rangle.
\]
In particular, $x$ and $z$ have infinite order.  This also follows directly
from the realization of $\HH$ by upper triangular integral $3\times3$
matrices with diagonal entries equal to one: take
$x=I+E_{12}$, $y=I+E_{23}$, and $z=I+E_{13}$.

Choose an orientation of each unoriented edge of a graph $X$, and associate
a distinct variable $t_e$ to each such edge.  An edge path determines a word
by reading $t_e$ on a positive traversal and $t_e^{-1}$ on a negative
traversal.  Its word is an \emph{$\HH$-law} if it evaluates to the identity
for every independent assignment $t_e\in\HH$.  Assign each unoriented edge
a positive length $\ell(e)$, and count every traversal with this length.
For a subgraph $S$, set $\ell(S)=\sum_{e\in E(S)}\ell(e)$.
Our application has unit edge lengths.  The weighted formulation also
allows lengths to record subdivision; alternatively one can apply the
Lemma directly to the subdivided graph with a separate variable on each
unit edge.

The edge-neutral analogue, in which each edge is traversed equally
often in both directions, is the bicycle characterization of abelian
girth in Friedman--Izsak--Silberman \cite[Lemma~2.6]{FIS15}.
The two-variable Heisenberg substitutions below give the stronger
conclusion that every edge of a suitable bicycle is traversed at least
four times.

\begin{lemma}[Heisenberg laws contain a bicycle used four times]
\label{lem:heisenberg-law}
Let $W$ be a nonempty closed edge path in $X$, reduced as a based edge
path: no two consecutive traversals in its linear ordering are mutually
inverse.  Suppose the word of $W$ is a $\HH$-law.  Then its support contains
a bicycle $B$ such that every edge of $B$ is traversed at least four times
by $W$.  Consequently,
\[
  \ell(W)\geq 4\ell(B).
\]
In particular, if every bicycle in the support of $W$ has length at least
$T$, then $\ell(W)\geq4T$.
\end{lemma}

\begin{proof}
We first use one-variable substitutions to balance every edge.  We
then use two-variable substitutions to isolate a balanced closed path
supported on edges traversed at least four times.  Its support will
contain the required bicycle.

For an oriented edge $e$, let $n_e(V)$ denote the number of its positive
traversals minus the number of its negative traversals in an edge word
$V$.  Assigning $t_e=x$ and all other variables the identity gives
$x^{n_e(W)}=1$, so
\begin{equation}
\label{eq:heisenberg-edge-balance}
  n_e(W)=0\qquad\text{for every edge }e.
\end{equation}
Thus every used edge has positive even traversal multiplicity.  Call an
edge \emph{paired} if its multiplicity in $W$ is exactly two, and call it
\emph{heavy} if its multiplicity is at least four.

We first examine a paired edge $e$.  Its two signs are opposite by
\eqref{eq:heisenberg-edge-balance}.  Starting the cyclic word at its first
occurrence, write
\[
  W=t_e^{\varepsilon}U t_e^{-\varepsilon}V,
  \qquad \varepsilon\in\{1,-1\}.
\]
Cyclically rotating a law preserves lawhood, since a cyclic rotation of
a word is a conjugate of it under every assignment.  Fix an edge $f\ne e$
and assign $t_e=x$, $t_f=y$, and all remaining variables the identity.
Put $n=n_f(U)$.  Equation~\eqref{eq:heisenberg-edge-balance} gives
$n_f(V)=-n$, and the assigned word becomes
\[
  x^{\varepsilon}y^n x^{-\varepsilon}y^{-n}
     =z^{\varepsilon n}.
\]
The infinite order of $z$ implies $n_f(U)=0$.  The edge $e$ does not occur
inside $U$, so we have proved the following interval property:
\begin{equation}
\label{eq:heisenberg-paired-interval}
  n_f(U)=0\qquad\text{for every edge }f.
\end{equation}
Thus the open interval between the occurrences of a paired edge is
balanced on every edge.  The two-variable substitution gives the
required pairwise area condition directly.

Fix the original linear ordering of the occurrences in $W$, and join the
two positions of each paired edge by an interval.  These intervals
cannot cross.  Indeed, if the intervals for $e$ and $f$ crossed, the open
interval for $e$ would contain exactly one occurrence of $f$, contradicting
\eqref{eq:heisenberg-paired-interval}.  If paired edges exist, choose an
inclusion-minimal paired interval, and let $U$ be its interior edge path.
Every edge occurring in $U$ is heavy: if a paired edge occurred there,
the interval property would force both of its occurrences to lie there,
giving a smaller paired interval.

Moreover, $U$ is a closed edge path, because the inside endpoints of the
two opposite traversals of the same edge coincide.  It is nonempty,
since an empty interior would give an immediate reversal in $W$.  It is
reduced as a based path, being a contiguous subpath of $W$, and it is
balanced by \eqref{eq:heisenberg-paired-interval}.  If there are no paired
edges, take $U=W$; these same properties follow from the hypotheses and
\eqref{eq:heisenberg-edge-balance}, and every edge of $U$ is again heavy
in $W$.

We claim that the connected support $S$ of $U$ has first Betti number at
least two.  Here it matters that $U$ need not be cyclically reduced at
its basepoint.  Repeatedly remove its first and last traversals whenever
they are mutually inverse.  This expresses $U$ as $P C P^{-1}$, where
$C$ is a cyclically reduced closed path.  The path $C$ is nonempty: if this
process reached the empty path, its last removal would have removed two
adjacent inverse traversals from the original reduced path.  Each
removal subtracts opposite traversals of the same edge, so $C$ is still
balanced.

Deleting a vertex of valence one and its incident edge from a finite
graph cannot affect a cyclically reduced closed path: a visit to that
vertex would require an immediate reversal.  Such deletions also
preserve first Betti number as long as an edge remains.  If $S$ had
first Betti number zero, this deletion process would leave no edges,
contradicting the existence of $C$.  If $S$ had first Betti number one,
the remaining nonempty core would be a single cycle.  To see the last
assertion, the core is connected with minimum valence two and satisfies
$|E|=|V|$; hence all its valences are exactly two.  A nonempty cyclically
reduced path on a cycle traverses it a nonzero number of times in one
fixed direction.  Its exponent sum on each cycle edge is therefore
nonzero, contradicting the balance of $C$.  This proves the claim,
including when the basepoint of $U$ lies on a path leading to its cyclic
core.

Finally, choose a spanning tree of $S$ and two distinct edges outside
that tree.  Let $T'$ be the smallest subtree containing the endpoints
of those two edges, and let $B$ be their union with $T'$.  The subgraph
$B$ is connected and has first Betti number two.  Every leaf of $T'$
is an endpoint of one of the added edges, so $B$ has no vertices of
valence one; the case where $T'$ is a single vertex gives the same
conclusion directly.  Thus $B$ is a bicycle contained in $S$.  Each
of its edges is heavy in $W$, and positivity of the edge lengths gives
\[
  \ell(W)
    =\sum_{e\in\supp(W)}
       \bigl(\text{traversal multiplicity of }e\bigr)\ell(e)
    \geq4\sum_{e\in E(S)}\ell(e)
    \geq4\ell(B).
\]
The asserted lower bound follows.
\end{proof}

At a fixed comparison scale $m>0$, exclusion of bicycles of length
strictly below $1.52m$ therefore gives
\[
  \ell(W)\geq4(1.52)m=6.08m
\]
for any $\HH$-law $W$ whose support lies in the region to which that
exclusion applies.  In particular, every nonempty cyclically reduced
closed path there of length less than $6.03m$ is a nonlaw.
The final proposition uses the gap between $6.08$ and $6.04$ to allow
scale variation along a path; the smaller cutoff $6.03$ leaves room
for the integer rounding in the voltage selection.

The factor four in the lemma is sharp, and the lower bound is
non-strict.  Indeed, on a bouquet with two edges
$a,b$, the cyclically reduced word
\[
  aba^{-1}b^{-1}a^{-1}bab^{-1}
      =[a,b][a^{-1},b]
\]
is a $\HH$-law and has length $4\ell(a)+4\ell(b)$, attaining equality
for its bicycle support.  The same equality holds after subdividing
the two cycles, including subdivisions with even cycle lengths.


\subsection{Simultaneous small deletion in every layer}

We now combine the two expectation bounds with the deterministic word
lemma.  The decay in $H_i$ makes the total expected deleted fraction
summable over all heights, so one table realization has every required
property at once.

\begin{proposition}[Pruned old graph]
\label{prop:pruning}
Given any $\varepsilon>0$, the allowed parameters and the
universal tables can be chosen so that the additional two
deletions in this Section remove less than an $\varepsilon$
fraction of each $\mathscr A_i$ and $\mathscr B_i$.
In the surviving induced old graph $X_{\mathrm{old}}$:
\begin{enumerate}
 \item the folded labeling, padded-prefix recovery of
 Lemma~\ref{lem:tower-recovery}, and global word uniqueness of
 Proposition~\ref{prop:tower-labels} remain valid;
 \item a nonempty cyclically reduced closed path based at height
 $i$ has length greater than $.93m_i$;
 \item every bicycle containing a vertex of height $i$ has
 length at least $1.52m_i$;
 \item every nonempty cyclically reduced closed path based at
 height $i$ with length at most $6.04m_i$ is not a law in
 the independent edge variables of $\HH$.
\end{enumerate}
The expansion conclusions for the undeleted incidences hold
simultaneously, and the degrees remain fixed along the tower.
\end{proposition}

\begin{proof}
For a realization, let $Z_i^A,Z_i^B$ be the additional deleted
fractions.  Bounds on rooted objects imply bounds on all their
vertices, with only polynomial and small scale factors, as
follows.  An object of length $\ell$ rooted at height $j$
can meet layer $i$ only if $|i-j|\leq\ell$.  Every counted
object satisfies $\ell\leq4H_j/s$.  From
$H_j=H_i+b(j-i)$ and $s\geq20b$ we obtain
\[
 |i-j|\leq5H_i/s,
 \qquad |H_j-H_i|\leq5bH_i/s.
\]
Its vertex count is at most $\ell+1\leq5H_i/s+1$.
There are at most $10H_i/s+1$ possible root layers and at
most a linear number of lengths.  Lemmas~\ref{lem:pruning-return-count}
and~\ref{lem:pruning-theta-count} therefore imply, after division
by $|\mathscr B_i|=2^{H_i}$ or
$|\mathscr A_i|=2^{H_i+2r-s}$,
\begin{equation}
\label{eq:pruning-deletion-expectation}
 \mathbb E(Z_i^A+Z_i^B)\leq2^{-.01H_i}
\end{equation}
for sufficiently large $s/b$ and $H_0/s$.  More explicitly,
the larger rooted bound contributes
$2^{-.025H_i+.975(5bH_i/s)}$ times a fixed polynomial in
$H_i$.  The coefficient remains below $-.02$ for small
$b/s$, and that polynomial is absorbed by another $.01H_i$
after increasing $H_0$.  This proves
\eqref{eq:pruning-deletion-expectation}, including the cost
of deleting vertices of an object whose qualifying root is
in a different layer.

Since $H_i=H_0+bi$,
\[
 \sum_{i\geq0}\mathbb E(Z_i^A+Z_i^B)
 \leq\frac{2^{-.01H_0}}{1-2^{-.01b}}.
\]
Choose $H_0$ still larger so the last expression is less
than $\varepsilon/2$.  Markov's inequality then gives
positive probability that
$\sum_i(Z_i^A+Z_i^B)<\varepsilon$.  On that event each
individual deleted fraction is below $\varepsilon$.
Every table realization in the sequential law already has
the undeleted expansion properties, so choose one in this
positive-probability event.  The earlier good-vertex deletion
budget is separately supplied by Proposition~\ref{prop:tower-labels}
and can be made as small as required there.

Passing to an induced subgraph preserves foldedness and
uniqueness of a word that occurred at most once before
deletion.  Lemma~\ref{lem:pruning-base-girth} proves the
second assertion, and the definition of the second deletion
proves the third.  It remains to deduce the nonlaw assertion, allowing
for the variation of scale along the closed path.

Let $W$ be a nonempty cyclically reduced closed path of length
$L\leq6.04m_i$ based at height $i$ in $X_{\mathrm{old}}$.
Every vertex $v$ of its support has scale
$ m(v)\geq m_i-aL\geq(1-6.04a)m_i$.
If $W$ were a $\HH$-law, Lemma~\ref{lem:heisenberg-law}
would give a bicycle $B$ in its support with $L\geq4|E(B)|$.
The third assertion, already proved by the deletion, gives
\[
 |E(B)|\geq1.52(1-6.04a)m_i.
\]
Consequently
\[
 L\geq6.08(1-6.04a)m_i>6.04m_i,
\]
where the last inequality holds for $a\leq10^{-6}$.
This contradicts the choice of $W$.  Independent edge variables
here include every underlying edge, regardless of any scalar
weight assigned later; reverse traversals use the inverse of
the same variable.
\end{proof}
